\chapter{Glossary}
\label{app:glossary}
\automark[section]{chapter}
\markboth{Glossary}{Glossary}

The definitions below explain the terms used in the thesis. A study's own
metric definition takes precedence where the same abbreviation has several
meanings.

\begin{description}
\item[Ablation] An experiment that removes or changes one part of a system
to estimate its contribution. Other relevant conditions should be held fixed.
\item[AC@k in CIRCA] The average fraction of labelled root-cause metrics
recovered among the first $k$ recommendations. It is not automatically the
percentage of incidents diagnosed completely correctly.
\item[Agent] A system that uses a model to choose actions, call tools and
use their results when deciding what to do next.
\item[Anomaly] An observation that differs from an expected pattern.
An anomaly can be a symptom rather than the cause of a failure.
\item[Apache Flink job] A data-processing task executed on the distributed
Apache Flink platform. An anomalous job shows unexpected behaviour.
\item[Baseline] A reference method against which an evaluated method is compared.
\item[Calibration] Agreement between a model's stated confidence and its
observed correctness across comparable cases.
\item[Confounding] A situation where several factors change together,
preventing a result from being attributed to one of them alone.
\item[Counterfactual] A hypothetical change used to ask what would have
happened if a suspected condition had been different.
\item[Evidence grounding] Linking a diagnostic claim to observations that
support it. A cited observation can still be misinterpreted.
\item[Expected calibration error (ECE)] A summary of the differences between
confidence and observed correctness across confidence groups. The grouping
and reference labels affect the result.
\item[False positive] A reported fault where the reference case is healthy
or does not contain that fault. The denominator must be stated.
\item[Ground truth] The reference used to judge correctness, such as a
recorded injected fault. Its quality must also be checked.
\item[LLM] A large language model, trained to model and generate language.
An LLM is itself a form of machine learning.
\item[Macro-F1] The unweighted average of the F1 scores calculated separately
for each class. Rare classes receive the same weight as common classes.
\item[Micro-F1] F1 calculated after combining correct predictions, false
positives and false negatives across classes. F1 balances precision and
recall. In single-label multiclass classification, micro-F1 equals accuracy.
\item[Mean reciprocal rank (MRR)] The average reciprocal of the first correct
answer's rank. A first-place answer contributes 1 and a second-place answer
contributes $1/2$. The treatment of missing answers must be specified.
\item[Multimodal telemetry] Different kinds of observations used together,
such as logs, metrics and traces.
\item[Precision and recall] Precision measures how many selected items are
correct. Recall measures how many relevant items were recovered.
\item[Preprint] A publicly available research manuscript whose retrieved
version is not necessarily peer reviewed.
\item[Retrieval-augmented generation (RAG)] Providing a language model with
selected external information before it generates an answer.
\item[Root-cause analysis (RCA)] Investigating the condition responsible for
a failure. Locating a faulty component is one part of that investigation.
\item[Span and trace] A span represents one operation. A trace groups related
spans that describe a request's execution across components.
\item[Telemetry] Observations emitted or collected while a system runs,
including logs, metrics, traces and events.
\item[Token] A unit of model input or output, often a word fragment. Token
counts help describe inference cost and context limits.
\item[Top-k localisation] Whether the correct component appears among the
first $k$ candidates. Multi-cause tasks require an explicit scoring rule.
\item[WAF] A web application firewall that applies rules to HTTP requests.
\end{description}
